% Generated draft: requires mathematical and source review.
\documentclass[11pt]{article}
\usepackage[margin=25mm]{geometry}
\usepackage{fontspec}
\setmainfont{Latin Modern Roman}
\usepackage{amsmath,amssymb,mathtools,unicode-math}
\setmathfont{Latin Modern Math}
\usepackage{enumitem,needspace,xurl,hyperref}
\setlength{\parindent}{0pt}
\setlength{\parskip}{5pt}
\setlength{\emergencystretch}{3em}
\newcommand{\N}{\mathbb{N}}
\newcommand{\Z}{\mathbb{Z}}
\newcommand{\Q}{\mathbb{Q}}
\newcommand{\R}{\mathbb{R}}
\newcommand{\C}{\mathbb{C}}
\newcommand{\F}{\mathbb{F}}
\newcommand{\A}{\mathbb{A}}
\newcommand{\PP}{\mathbb{P}}
\newcommand{\E}{\mathbb{E}}
\newcommand{\eps}{\varepsilon}
\newcommand{\dd}{\,\mathrm{d}}
\begin{document}
\section*{UnsolvedMath ID: 30003365}
\textbf{Source problem code:} catalogue entry
\subsection{Definitions and mathematical statement}
Let \((S,g)\) be a Riemannian surface of genus at least \(2\), with Gaussian curvature \(\kappa_g\leq 0\), and write
\[
T^1S=\{v\in TS:g(v,v)=1\}.
\]
For \(v\in T^1S\), let \(\gamma_v:I_v\to S\) be its maximal unit-speed geodesic, with \(\dot\gamma_v(0)=v\). Define its rank as the dimension of the space of parallel Jacobi fields along \(\gamma_v\). On a surface, rank \(2\) is equivalent to the existence of a nonzero parallel normal Jacobi field, and hence to \(\kappa_g(\gamma_v(t))=0\) for every \(t\in I_v\). Thus the set in the supplied problem is
\[
K=\{v\in T^1S:\operatorname{rank}(\gamma_v)=2\}
 =\{v\in T^1S:\kappa_g(\gamma_v(t))=0\text{ for all }t\in I_v\}.
\]
If the metric is geodesically complete, then \(I_v=\R\) for every \(v\).

\textbf{Requested conclusions.} Must \(K\) have measure zero, and must it consist of finitely many flat bands of closed geodesics? The natural measure interpretation is the Liouville measure \(\mu_L\) on \(T^1S\), locally given by Riemannian area on \(S\) times angular measure on the unit tangent circle. Under this explicitly provisional interpretation, the first question is
\[
\mu_L(K)=0\;?
\]
A flat band ordinarily refers to a Euclidean product cylinder \((\R/\ell\Z)\times J\), with metric \(\dd s^2+\dd r^2\), mapped isometrically into \(S\), whose curves with constant \(r\) are closed geodesics. Its corresponding subset of \(T^1S\) consists of the unit tangent vectors to those curves in both directions. With a specified convention for such bands, the second question asks whether
\[
K=\bigcup_{i=1}^{m}B_i
\qquad\text{for some finite }m\geq 0,
\]
where each \(B_i\) is the corresponding set of tangent vectors for one band.

\textbf{Unresolved qualifications.} The source does not specify the measure, compactness, completeness, orientability, boundary conditions, or metric regularity beyond the Riemannian curvature terminology. It also does not specify whether bands are embedded or immersed, which transverse intervals \(J\) are allowed, or whether isolated closed flat geodesics count as zero-width bands. These choices are not imposed here. In particular, the finite-union formulation remains dependent on the intended band convention; the supplied record does not determine a fully qualified conjecture.
\subsection{Short English statement}
For a nonpositively curved surface of genus at least \(2\), consider all unit tangent vectors whose geodesics remain entirely in zero Gaussian curvature. The problem asks whether this set has measure zero and is made up of finitely many flat bands of closed geodesics. The intended surface hypotheses, measure, and treatment of isolated closed flat geodesics require clarification.
\subsection{Sources}
\begin{itemize}
\item Local UnsolvedMath registry, permanent ID \(30003365\), supplied complete source record, title: \emph{Flat Bands on Nonpositively Curved Surfaces}; exact source problem code: \texttt{catalogue entry}. Collection: \emph{Oberwolfach Reports Open Problems}. No individual proposer is supplied; the proposed year is recorded as \(2017\).
\item Supplied report citation: \emph{Mini-Workshop: Spaces and Moduli Spaces of Riemannian Metrics (2018)}, Oberwolfach Reports, pp.\ 133--166, DOI \texttt{10.4171/owr/2017/3}. \url{https://doi.org/10.4171/owr/2017/3}. The record describes Oberwolfach Reports as documenting mathematical workshops and their open-problem discussions. The citation's displayed year \(2018\) differs from the year in the DOI and the recorded proposed year; this discrepancy is unresolved.
\item Supplied literature reference: Wu, \emph{On the ergodicity of geodesic flows}. \url{https://www.numdam.org/item/10.5802/afst.1457.pdf}.
\item Literature assessment checked at \texttt{2026-08-21}: supplied dated curation, not independently verified. The record labels the problem open and reports that the general finiteness conjecture remains open, while Wu proves a finite-negative-component case in which flat strips and isolated closed flat geodesics are finite. The record does not supply the precise hypotheses or theorem statement for that case. Its background expressly calls this dated literature triage derived from the supplied OWR curation and requests independent verification of time-sensitive status and cited claims.
\item The registry labels the statement \texttt{corrected\_verified}, explaining that notation, typography, and sentence structure were repaired without changing mathematical content. This supplied editorial label is not independent mathematical verification.
\end{itemize}
\end{document}
